% 原书物理页=印刷页25--31上部（2.5 标题之前）；2.4 全节. 2.4.4 的证明与 Cartan 判据在原书排至 p31 上部，仍属本节，一并译出.
\section{幂零与可解 Lie 代数}\label{sec:ssl-2-4}
设 $\mathfrak{g}$ 是一个 $K$-Lie 代数. 对 $\mathfrak{g}$ 的两个理想 $\mathfrak{h}_1, \mathfrak{h}_2$，把 $[\mathfrak{h}_1, \mathfrak{h}_2]$ 定义为由诸括号 $[x_1, x_2]$（$x_1 \in \mathfrak{h}_1$，$x_2 \in \mathfrak{h}_2$）张成的 $K$-向量子空间. 根据 Jacobi 恒等式，它仍是 $\mathfrak{g}$ 的理想.
\begin{example*}
（导理想）理想 $[\mathfrak{g}, \mathfrak{g}]$ 称为 $\mathfrak{g}$ 的\newterm{ssl-derived-ideal}{导理想}{derived ideal}. 它为零当且仅当 $\mathfrak{g}$ 是 Abel 的. 一般地，它含于 $\mathfrak{g}$ 到任一 Abel Lie 代数的每个态射之中\footnote{译注：原文如此（``含于每个态射之中''）；按文义应指含于该态射的核中.}，从而 $\mathfrak{g}_{\mathrm{ab}} := \mathfrak{g}/[\mathfrak{g}, \mathfrak{g}]$ 是 $\mathfrak{g}$ 的最大 Abel 商.
\end{example*}
更一般地，定义两个（按包含关系）递降的理想链：
\begin{enumerate}[label=\roman*)]
\item \newterm{ssl-lower-central-series}{降中心列}{lower central series}，定义为 $C^0\mathfrak{g} := \mathfrak{g}$ 且 $C^{n+1}\mathfrak{g} := [\mathfrak{g}, C^n\mathfrak{g}]$（$n \geqslant 0$）；
\item \newterm{ssl-derived-series}{导列}{derived series}，定义为 $D^0\mathfrak{g} := \mathfrak{g}$ 且 $D^{n+1}\mathfrak{g} := [D^n\mathfrak{g}, D^n\mathfrak{g}]$（$n \geqslant 0$）.
\end{enumerate}
有 $C^1\mathfrak{g} = D^1\mathfrak{g} = [\mathfrak{g}, \mathfrak{g}]$，且对 $n \geqslant 1$ 有 $C^n\mathfrak{g} \supset D^n\mathfrak{g}$.
\begin{definition}\label{def:ssl-2-4-1}
称 $\mathfrak{g}$ 是
\begin{itemize}[label=\textendash]
\item \newterm{ssl-nilpotent}{幂零}{nilpotent}的，如果存在 $n \in \mathbb{N}$ 使 $C^n\mathfrak{g} = 0$；
\item \newterm{ssl-solvable}{可解}{solvable}的，如果存在 $n \in \mathbb{N}$ 使 $D^n\mathfrak{g} = 0$.
\end{itemize}
于是有 Abel $\Rightarrow$ 幂零 $\Rightarrow$ 可解. 还应注意到这里与有限群的幂零性、可解性概念之间的平行.
\end{definition}
\begin{exercise*}
设 $\mathfrak{g}$ 是一个 $K$-Lie 代数.
\begin{enumerate}[label=\roman*)]
\item
\begin{enumerate}[label=(\alph*)]
\item 证明 $\mathfrak{g}$ 幂零当且仅当存在 $\mathfrak{g}$ 的理想链 $\mathfrak{g} = \mathfrak{h}_0 \supset \mathfrak{h}_2 \supset \cdots \supset \mathfrak{h}_n = 0$，使得对一切 $i$ 有 $[\mathfrak{g}, \mathfrak{h}_i] \subset \mathfrak{h}_{i+1}$.\footnote{译注：链中第二项的下标原文两处（i)(a) 与 ii)(a)）均印作 $2$，按文义应为 $\mathfrak{h}_1$.}
\item 由此推出：幂零 Lie 代数的任何 Lie 子代数与任何商代数都是幂零的.
\item 证明 $\mathfrak{g}$ 幂零当且仅当 $\operatorname{ad}(\mathfrak{g}) = \mathfrak{g}/\mathfrak{z}(\mathfrak{g})$ 幂零.
\end{enumerate}
\item
\begin{enumerate}[label=(\alph*)]
\item 证明 $\mathfrak{g}$ 可解当且仅当存在 $\mathfrak{g}$ 的理想链 $\mathfrak{g} = \mathfrak{h}_0 \supset \mathfrak{h}_2 \supset \cdots \supset \mathfrak{h}_n = 0$，使得对一切 $i$ 有 $\mathfrak{h}_i/\mathfrak{h}_{i+1}$ 是 Abel 的.
\item 由此推出：若 $\mathfrak{h}$ 是 $\mathfrak{g}$ 的理想，则 $\mathfrak{g}$ 可解当且仅当\footnote{注意：把``可解''换成``幂零''后，此处的``若''方向不成立.} $\mathfrak{h}$ 与 $\mathfrak{g}/\mathfrak{h}$ 都可解.
\end{enumerate}
\end{enumerate}
\end{exercise*}
\begin{definition}\label{def:ssl-2-4-2}
设 $V$ 是一个 $n$ 维 $K$-向量空间. $V$ 的一个\newterm{ssl-complete-flag}{完全旗}{complete flag}是指子空间的一个递增序列 $\mathcal{V} = (0 = \mathcal{V}_0 \subset \mathcal{V}_1 \subset \cdots \subset \mathcal{V}_n = V)$，其中 $\dim(\mathcal{V}_i) = i$. 与之相伴有 $\mathfrak{gl}(V)$ 的两个 Lie 子代数
\[\begin{gathered}
\mathfrak{b}(\mathcal{V}) := \{X \in \mathfrak{gl}(V),\ X(\mathcal{V}_i) \subset \mathcal{V}_i,\ i = 1, \cdots, n\},\ \text{以及}\\
\mathfrak{b}^{\mathrm{nilp}}(\mathcal{V}) := \{X \in \mathfrak{gl}(V),\ X(\mathcal{V}_i) \subset \mathcal{V}_{i-1},\ i = 1, \cdots, n\}.
\end{gathered}\]
\end{definition}
\begin{remark*}
若选取一组基 $v_1, \cdots, v_n$ 使 $\mathcal{V}_i = \operatorname{vect}(v_1, \cdots, v_i)$，则由此得到的同构 $\mathfrak{gl}(V) \xrightarrow{\ \sim\ } \mathfrak{gl}_n(K)$ 便把 $\mathfrak{b}(\mathcal{V})$，相应地 $\mathfrak{b}^{\mathrm{nilp}}(\mathcal{V})$，等同于 $M_n(K)$ 中的上三角矩阵空间，相应地严格上三角矩阵空间.
\end{remark*}
\begin{exercise*}
验证 $\mathfrak{b}(\mathcal{V})$ 确是 $\mathfrak{gl}(V)$ 的 Lie 子代数，然后证明：
\begin{enumerate}[label=\roman*)]
\item $[\mathfrak{b}(\mathcal{V}), \mathfrak{b}(\mathcal{V})] = \mathfrak{b}^{\mathrm{nilp}}(\mathcal{V})$，且 $[\mathfrak{b}(\mathcal{V}), \mathfrak{b}^{\mathrm{nilp}}(\mathcal{V})] = \mathfrak{b}^{\mathrm{nilp}}(\mathcal{V})$；
\item 对 $k \geqslant 1$，$C^k(\mathfrak{b}^{\mathrm{nilp}}(\mathcal{V})) = \{X \in \mathfrak{gl}(V),\ X(\mathcal{V}_i) \subset \mathcal{V}_{i-k},\ i = 1, \cdots, n\}$.\footnote{译注：此式中的 $i-k$ 原文如此. 按前文 $C^0\mathfrak g=\mathfrak g$ 的约定，严格上三角矩阵本身已使旗至少下降一级，因此此处应为 $i-k-1$；原式照录.}
\end{enumerate}
由此推出 $\mathfrak{b}^{\mathrm{nilp}}(\mathcal{V})$ 是幂零的，而 $\mathfrak{b}(\mathcal{V})$ 是可解的但并非幂零的.
\end{exercise*}

\begin{numpara}[借助表示论的刻画]\label{np:ssl-2-4-3}
这里假设 $K$ 是特征为零的代数闭域，而且不妨只考虑 $K = \mathbb{C}$.
我们将首先看到，上面那个可解 Lie 代数的例子是典型的.
\begin{theorem*}[Lie]
设 $\mathfrak{g}$ 是一个 $K$-Lie 代数，$K$ 代数闭. 下列条件等价：
\begin{enumerate}[label=\roman*)]
\item $\mathfrak{g}$ 可解；
\item $\mathfrak{g}$ 的每个不可约表示都是一维的；
\item[ii)$'$] $\mathfrak{g}$ 的每个表示都含一个对 $\mathfrak{g}$ 中所有元素公共的特征向量；
\item $\mathfrak{g}$ 的每个表示 $(V, r)$ 都可\newterm{ssl-trigonalizable}{上三角化}{trigonalizable}，意思是存在 $V$ 的一个旗 $\mathcal{V}$，使 $r(\mathfrak{g}) \subset \mathfrak{b}(\mathcal{V})$.
\end{enumerate}
\end{theorem*}
\begin{proof}
等价 ii)$\Leftrightarrow$ii)$'$ 是明显的，蕴涵 iii)$\Rightarrow$ii)$'$ 同样明显，而蕴涵 ii)$'\Rightarrow$iii) 由一个显然的归纳即可看出.
往证 iii)$\Rightarrow$i). 把 iii) 应用于伴随表示. 我们得到子空间的一个完全旗 $0 \subset \mathfrak{h}_1 \subset \cdots \subset \mathfrak{h}_n = \mathfrak{g}$，使得 $\operatorname{ad}_{\mathfrak{g}}(\mathfrak{h}_i) \subset \mathfrak{h}_i$，这就意味着诸 $\mathfrak{h}_i$ 是 $\mathfrak{g}$ 的理想. 由于商 $\mathfrak{h}_i/\mathfrak{h}_{i+1}$ 是一维的\footnote{译注：原文商的次序如此. 这里的旗递增，因此应为 $\mathfrak h_{i+1}/\mathfrak h_i$，原式照录.}，它是一个 Abel Lie 代数，于是可以应用上面习题中的判别法，推出 $\mathfrak{g}$ 可解.
现在对 $\dim(\mathfrak{g})$ 作归纳，往证 i)$\Rightarrow$ii)$'$. 若 $\dim(\mathfrak{g}) = 1$，结论是清楚的，因为有限维 $K$-向量空间上的任何自同态都有特征向量（$K$ 代数闭）. 设 $\dim(\mathfrak{g}) > 1$，取 $\mathfrak{g}$ 中包含 $D\mathfrak{g}$ 的一个超平面 $\mathfrak{h}$（注意 $\mathfrak{g}/D\mathfrak{g} \neq 0$）. 那么 $\mathfrak{h}$ 是 $\mathfrak{g}$ 的理想，且由归纳假设，存在 $V$ 的一个向量 $v$，它同时是 $\mathfrak{h}$ 中所有元素的特征向量. 设 $\lambda$ 是 $\mathfrak{h}$ 上由 $X.v = \lambda(X)v$（对一切 $X \in \mathfrak{h}$）定义的线性型. 置
\[V_{\lambda} := \{v \in V,\ \forall X \in \mathfrak{h},\ X.v = \lambda(X)v\},\]
即 $V$ 在 $\mathfrak{h}$ 的作用下的``$\lambda$-特征''子空间，它是非零的. 暂且承认 $V_{\lambda}$ 在 $\mathfrak{g}$ 下稳定. 取 $\mathfrak{h}$ 在 $\mathfrak{g}$ 中一条补直线的生成元 $Y$. 由于 $K$ 代数闭，$Y$ 在 $V_{\lambda}$ 中有一个特征向量 $v'$. 于是 $v'$——它已是 $\mathfrak{h}$ 的特征向量——便是整个 $\mathfrak{g}$ 的特征向量，定理得证.
剩下要证 $V_{\lambda}$ 在 $\mathfrak{g}$ 下稳定. 由等式
\[XYv = YXv + [X, Y]v = \lambda(X)Yv + \lambda([X, Y])v\]
（它对一切 $X \in \mathfrak{h}$，$Y \in \mathfrak{g}$，$v \in V_{\lambda}$ 成立），我们被归结为证明：对一切 $X \in \mathfrak{h}$，$Y \in \mathfrak{g}$，有 $\lambda([X, Y]) = 0$. 为此，固定一个向量 $w \neq 0$（$w \in V_{\lambda}$），并记
\[W_k := \operatorname{vect}\{w, Yw, \cdots, Y^k w\}.\]
公式 $XY^k w = YXY^{k-1}w + [X, Y]Y^{k-1}w$ 由归纳表明
\[\forall X \in \mathfrak{h},\quad XW_k \subset W_k\quad\text{且}\quad (\ast):\ XY^k w \in \lambda(X)Y^k w + W_{k-1}.\]
于是考虑诸 $W_k$ 的并 $W$. 它是一个 $\mathfrak{h}$-稳定的子空间，且 $(\ast)$ 表明对一切 $X \in \mathfrak{h}$ 有 $\operatorname{tr}(X|_W) = \dim(W)\lambda(X)$. 特别地，$\dim(W)\lambda([X, Y]) = \operatorname{tr}([X, Y]|_W)$. 而按构造，$W$ 在 $Y$ 下也稳定，故 $\operatorname{tr}([X, Y]|_W) = \operatorname{tr}([X|_W, Y|_W] = 0$\footnote{译注：原文此式缺右括号，照录；应为 $\operatorname{tr}([X|_W, Y|_W]) = 0$.}. 于是 $\lambda([X, Y]) = 0$，正如所需（因为在 $K$ 特征为零的假设下，$\dim(W)$ 在 $K$ 中非零）.
\end{proof}
\begin{remark*}
本定理推广了线性代数中的经典结果：两两交换的自同态族可以同时上三角化.
\end{remark*}
为了通过表示来刻画幂零 Lie 代数，我们引入几个定义与记号. 设 $\lambda : \mathfrak{g} \longrightarrow K$ 是一个映射，$(V, r)$ 是一个表示，记
\[V_{\lambda} := \{v \in V,\ \forall X \in \mathfrak{g},\ r(X)v = \lambda(X)v\}\]
为相伴的特征空间，而
\[V^{\lambda} := \{v \in V,\ \exists n \in \mathbb{N},\ \forall X \in \mathfrak{g},\ (r(X) - \lambda(X))^n v = 0\}\]
为相伴的\newterm{ssl-generalized-eigenspace}{广义特征空间}{generalized eigenspace}（或称``特征''空间）.
\begin{definition*}
若 $V_{\lambda} \neq 0$，则称 $\lambda$ 是 $\mathfrak{g}$ 在表示 $(V, r)$ 中的一个权. 这时 $V_{\lambda}$ 称为权 $\lambda$ 的\newterm{ssl-weight-space}{权空间}{weight space}，而 $V^{\lambda}$ 称为\newterm{ssl-generalized-weight-space}{广义权空间}{generalized weight space}.
\end{definition*}
\begin{exercise*}
若 $\lambda$ 是某个表示的权，则 $\lambda$ 是 $K$-线性的，且 $\lambda([\mathfrak{g}, \mathfrak{g}]) = 0$. 因此 $\lambda$ 是 $\mathfrak{g}$ 的\newterm{ssl-abelianization}{Abel 化}{abelianization} $\mathfrak{g}_{\mathrm{ab}}$ 的对偶 $\mathfrak{g}^{*}_{\mathrm{ab}}$ 中的元素.
\end{exercise*}
注意，一般说来，$V^{\lambda}$ 是否在 $\mathfrak{g}$ 下稳定并不清楚.
\begin{proposition*}[Engel 定理]
若 $V = V^{\lambda}$，则 $\lambda \in \mathfrak{g}^{*}_{\mathrm{ab}}$，且存在 $V$ 的完全旗 $\mathcal{V}$，使 $(r - \lambda)(\mathfrak{g}) \subset \mathfrak{b}^{\mathrm{nilp}}(\mathcal{V})$.
\end{proposition*}
\begin{remark*}
具体地说，这个命题断言：如果对每个 $X$ 都存在一组基，使 $r(X)$ 在这组基下是对角线上取 $\lambda(X)$ 的上三角矩阵，那么事实上存在一组基，使所有 $r(X)$ 都具有这种形状.
\end{remark*}
\begin{proof}
按假设，对一切 $X \in \mathfrak{g}$ 有 $\operatorname{tr}(r(X)) = \lambda(X)\dim(V)$. 这表明（设 $K$ 特征为 $0$，以使 $\dim(V)$ 在 $K$ 中非零）$\lambda$ 是 $K$-线性的且在 $[\mathfrak{g}, \mathfrak{g}]$ 上为零，故它来自 $\mathfrak{g}^{*}_{\mathrm{ab}}$. 于是映射 $X \longmapsto r(X) - \lambda(X)\,\mathrm{id}_V$ 在同一个空间 $V$ 上定义了 $\mathfrak{g}$ 的一个表示. 用 $r - \lambda$ 替换 $r$，便可设 $\lambda = 0$. 由一个递归推理，这时只须证明：
\[(\ast)\quad\text{对每个满足``对一切 } X \text{ 有 } r(X) \text{ 幂零''的表示 } (V, r)\text{，都有 } V^{\mathfrak{g}} \neq 0.\]
我们对 $\dim(\mathfrak{g})$ 作归纳来证明它. 若 $\dim(\mathfrak{g}) = 1$，则对 $\mathfrak{g}$ 的每个非零元素 $X$ 都有 $V^{\mathfrak{g}} = \operatorname{Ker}(r(X))$. 既然 $r(X)$ 是幂零的，它的核非零，得证.
设 $\dim(\mathfrak{g}) > 1$，并假设我们有一个非平凡理想 $\mathfrak{h}$，即 $\mathfrak{h}$ 既不等于 $\{0\}$ 也不等于 $\mathfrak{g}$. 那么把归纳假设应用于 $\mathfrak{h}$，便得 $V^{\mathfrak{h}} \neq \{0\}$. 此外，$V^{\mathfrak{h}}$ 在 $\mathfrak{g}$ 的作用下稳定，且该作用通过 $\mathfrak{g}/\mathfrak{h}$ 分解. 于是把归纳假设应用于 $\mathfrak{g}/\mathfrak{h}$，便得 $(V^{\mathfrak{h}})^{\mathfrak{g}/\mathfrak{h}} \neq 0$. 而 $V^{\mathfrak{g}} = (V^{\mathfrak{h}})^{\mathfrak{g}/\mathfrak{h}}$.
于是只剩下证明非平凡理想 $\mathfrak{h}$ 的存在. 用 $r(\mathfrak{g})$ 替换 $\mathfrak{g}$ 之后，可设 $\mathfrak{g} \subset \mathfrak{gl}(V)$. 选取 $\mathfrak{g}$ 的一个维数最大的真 Lie 子代数 $\mathfrak{h}$. $\mathfrak{g}$ 的伴随表示 $\operatorname{ad}$ 限制在 $\mathfrak{h}$ 上，诱导出 $\mathfrak{h}$ 在商空间 $\mathfrak{g}/\mathfrak{h}$ 上的一个表示. $\mathfrak{h}$ 在该表示下的像由幂零自同态组成. 事实上，根据下面的习题，对每个 $X \in \mathfrak{g}$，$\mathfrak{g}$ 的自同态 $\operatorname{ad}(X)$ 都是幂零的. 于是可以应用我们的归纳假设，推出 $\mathfrak{g}/\mathfrak{h}$ 中存在被 $\mathfrak{h}$ 的作用零化的元素 $\bar{X} \neq 0$. 取 $\mathfrak{g}$ 中位于 $\bar{X}$ 之上的元素 $X$. 于是 $X \notin \mathfrak{h}$，且 $[\mathfrak{h}, X] = \operatorname{ad}(\mathfrak{h})(X) \subseteq \mathfrak{h}$. 这蕴涵 $\mathfrak{h} \oplus K \cdot X$ 是一个以 $\mathfrak{h}$ 为理想的 Lie 代数. 但由 $\mathfrak{h}$ 的最大性，有 $\mathfrak{h} \oplus K \cdot X = \mathfrak{g}$. 这样，$\mathfrak{h}$ 就是 $\mathfrak{g}$ 的余维数为 $1$ 的理想.
\end{proof}
\begin{exercise*}
设 $A$ 是一个结合 $K$-代数，对 $x, y \in A$ 令 $\operatorname{ad}(x)(y) = [x, y] := xy - yx$. 证明 $\operatorname{ad}(x)^n(y) \in \operatorname{Vect}_K\{x^k y x^{n-k},\ k = 0, \cdots, n\}$.
\end{exercise*}
\begin{corollary*}
设 $\mathfrak{g}$ 是一个有限维 $K$-Lie 代数. 那么 $\mathfrak{g}$ 幂零当且仅当对每个 $X \in \mathfrak{g}$，$\operatorname{ad}_X$ 都是幂零的.
\end{corollary*}
\begin{proof}
由于 $\operatorname{ad}^n_X(\mathfrak{g}) \subset C^n\mathfrak{g}$，可见该条件是必要的. 反之，设对每个 $X \in \mathfrak{g}$，$\operatorname{ad}_X$ 幂零. 根据上面的命题（取 $\lambda = 0$），存在一个旗 $\mathfrak{h}_1 \subset \mathfrak{h}_2 \subset \cdots \subset \mathfrak{h}_n = \mathfrak{h}$，使得对一切 $X$ 有 $\operatorname{ad}_X(\mathfrak{h}_i) \subset \mathfrak{h}_{i-1}$，亦即 $[\mathfrak{g}, \mathfrak{h}_i] \subset \mathfrak{h}_{i-1}$\footnote{译注：原文旗的末项作 $\mathfrak{h}_n = \mathfrak{h}$；按上下文（对伴随表示应用上命题）应为 $\mathfrak{h}_n = \mathfrak{g}$.}. 把 $\mathfrak{h}_i$ 重新编号为 $\mathfrak{h}_{n-i}$，就认出了定义之后的那个幂零性判别法.
\end{proof}
\begin{lemma*}
设 $\mathfrak{g}$ 是一个幂零 $K$-Lie 代数，$\mathfrak{h} \subset \mathfrak{g}$ 是一个 Lie 子代数，$\lambda : \mathfrak{h} \longrightarrow K$ 是一个映射，而 $V^{\lambda}$ 是广义特征空间. 那么 $V^{\lambda}$ 在 $\mathfrak{g}$ 下稳定.
\end{lemma*}
\begin{proof}
设 $Y \in \mathfrak{g}$，$X \in \mathfrak{h}$，$v \in V^{\lambda}$. 我们要证明
\[(r(X) - \lambda(X))^n r(Y)v = 0\quad\text{对足够大的 } n\ \text{成立}.\]
而对 $V$ 的两个自同态 $f, g$，可以用归纳法（习题）证明下述公式\footnote{译注：式中二项式系数原文印作上 $k$ 下 $n$（即 $\binom{k}{n}$），按数学应为 $\binom{n}{k}$，疑上下排倒.}：
\[f^n \circ g = \sum_{k=0}^{n} \binom{k}{n} (\operatorname{ad}_f)^k(g) \circ f^{n-k}.\]
置 $f = r(X) - \lambda(X)$，$g = r(Y)$. 取 $m$ 使 $f^m(v) = 0$ 且 $\operatorname{ad}^m_f(g) = 0$（事实上 $\operatorname{ad}_f = \operatorname{ad}_{r(X)}$，而 $\mathfrak{g}$ 幂零，故 $r(\mathfrak{g})$ 幂零）. 于是上面的公式给出 $f^{2m} \circ g(v) = 0$. 这样就证明了 $V^{\lambda}$ 在 $\mathfrak{g}$ 下稳定.
\end{proof}
\begin{theorem*}
设 $K$ 代数闭，$\mathfrak{g}$ 是一个有限维 $K$-Lie 代数. 则下列条件等价：
\begin{enumerate}[label=\roman*)]
\item $\mathfrak{g}$ 幂零；
\item 每个表示 $(V, r)$ 都有分解 $V = \bigoplus_{\lambda} V^{\lambda}$，其中 $\lambda$ 取遍 $\mathfrak{g}$ 在 $V$ 中的权，且 $V^{\lambda}$ 在 $\mathfrak{g}$ 下稳定.
\end{enumerate}
\end{theorem*}
\begin{proof}
ii)$\Rightarrow$i). 把 ii) 应用于伴随表示. 由此得到直和分解 $\mathfrak{g} = \bigoplus_{\lambda} \mathfrak{g}^{\lambda}$，其各项都是理想. 证明 $\mathfrak{g}^{\lambda} \neq 0 \Rightarrow \lambda = 0$. 设 $Y \in \mathfrak{g}^{\lambda}$. 自同态 $(\operatorname{ad}_Y - \lambda(Y))|_{\mathfrak{g}^{\lambda}}$ 是幂零的，而 $(\operatorname{ad}_Y)|_{\mathfrak{g}^{\lambda}}$ 不可逆，因为 $\operatorname{ad}_Y(Y) = 0$. 由此得 $\lambda(Y) = 0$，于是 $\lambda|_{\mathfrak{g}^{\lambda}} = 0$. 又因 $[\mathfrak{g}^{\lambda'}, \mathfrak{g}^{\lambda}] \subset \mathfrak{g}^{\lambda'} \cap \mathfrak{g}^{\lambda}$，对其余的权也有 $\lambda|_{\mathfrak{g}^{\lambda'}} = 0$，最终得 $\lambda = 0$. 于是由上面的推论，$\mathfrak{g}$ 是幂零的.
i)$\Rightarrow$ii). 对 $\dim(V)$ 作归纳. 若所有 $r(X)$ 都只有一个特征值，则对某个函数 $\lambda$ 有 $V = V^{\lambda}$，而我们已看到 $\lambda$ 是一个权（上面的命题）. 否则，存在 $X$ 至少有两个不同的特征值（$K$ 代数闭）. 由此得到非平凡分解
\[V = V^{\lambda_1}_X \oplus \cdots \oplus V^{\lambda_k}_X,\]
为 $r(X)$ 的广义特征子空间之直和. 由前一引理，每个 $V^{\lambda_i}_X$ 都在 $\mathfrak{g}$ 下稳定，故可对它应用归纳假设.
\end{proof}
\begin{remark*}
这个定理为 $r(\mathfrak{g})$ 中的元素给出了``同时 Jordan 标准形''.
\end{remark*}
\end{numpara}

\begin{numpara}[借助 Killing 型的刻画]\label{np:ssl-2-4-4}
仍设 $K$ 是特征为零的代数闭域.
\begin{proposition*}
设 $\mathfrak{g}$ 是 $\mathfrak{gl}_n(K)$ 的一个 $K$-Lie 子代数，满足对一切 $X, Y \in \mathfrak{g}$ 有 $\operatorname{tr}(XY) = 0$. 那么 $[\mathfrak{g}, \mathfrak{g}]$ 是幂零的，从而 $\mathfrak{g}$ 可解.
\end{proposition*}
\begin{remark*}
$\mathfrak{g}$ 的可解性由定义 2.3.1 之下习题的 ii)(b)\footnote{译注：原文引作``定义 2.3.1''；该习题实在定义~\ref{def:ssl-2-4-1} 之下，2.3.1 疑为 2.4.1 之误.}应用于 $\mathfrak{h} = [\mathfrak{g}, \mathfrak{g}]$ 与 Abel 的 $\mathfrak{g}/\mathfrak{h} = \mathfrak{g}_{\mathrm{ab}}$ 而得出.
\end{remark*}
\begin{proof}
为简单起见，我们设 $K = \mathbb{C}$ 或 $K = \bar{\mathbb{Q}} \subset \mathbb{C}$（注意，若 $K$ 的基数至多为连续统的势（at most second-countable），则已知（在承认不可数选择公理的前提下）$K$ 确实同构于 $\mathbb{C}$ 或 \mbox{$\bar{\mathbb{Q}} \subset \mathbb{C}$）.}\footnote{译注：本括号中的基数与同构断言按原文照录. 例如 $\overline{\mathbb Q(t)}$ 是可数的特征零代数闭域，超越次数为 $1$，既不同构于 $\overline{\mathbb Q}$，也不同构于 $\mathbb C$. 所述基数条件可用于把该域嵌入 $\mathbb C$，不能推出只有上述两种同构类型. 原文英文括注 ``at most second-countable'' 也不是这一基数条件的标准说法.} 由\newterm{ssl-engel-theorem}{Engel 定理}{Engel's theorem}，只须证明每个元素 $X \in [\mathfrak{g}, \mathfrak{g}] \subset \mathfrak{gl}_n(K)$ 都幂零，换言之，其特征值 $\lambda_1, \cdots, \lambda_n$ 全为零. 取多项式 $f \in K[T]$，使对一切 $i$ 有 $f(\lambda_i) = \bar{\lambda}_i$（例如取一个 Lagrange 插值多项式）. 设 $X = X_s + X_n$ 为 $X$ 的 Jordan 分解. 那么
\[\operatorname{tr}(f(X_s)X) = \sum_i \bar{\lambda}_i \lambda_i = \sum_i |\lambda_i|^2.\]
因此只须证明 $\operatorname{tr}(f(X_s)X) = 0$.
既然 $X \in [\mathfrak{g}, \mathfrak{g}]$ 是诸 $[Y, Z]$（$Y, Z \in \mathfrak{g}$）的线性组合，且
\[\operatorname{tr}(f(X_s)[Y, Z]) = \operatorname{tr}([f(X_s), Y]Z),\]
只须证明 $\operatorname{ad}_{f(X_s)}(\mathfrak{g}) \subset \mathfrak{g}$，以便应用假设. 而 $\mathfrak{gl}_n(K)$ 的自同态 $\operatorname{ad}_{f(X_s)}$ 与自同态 $\operatorname{ad}_{X_s}$ 可在同一组基下对角化. 更精确地说，若 $(e_1, \cdots, e_n)$ 是使 $X_s$ 对角化的 $K^n$ 的一组基（特征值为 $\lambda_1, \cdots, \lambda_n$），而 $(E_{ij})_{1 \leqslant i,j \leqslant n}$ 是 $\mathfrak{gl}_n(K)$ 的相应基，则 $E_{ij}$ 是 $\operatorname{ad}_{X_s}$ 的、相应地 $\operatorname{ad}_{f(X_s)}$ 的特征向量，特征值分别为 $\lambda_i - \lambda_j$ 与 $f(\lambda_i) - f(\lambda_j) = \overline{\lambda_i - \lambda_j}$. 由此，对任何满足 $g(\lambda_i - \lambda_j) = \overline{\lambda_i - \lambda_j}$ 的多项式 $g \in K[T]$（这样的多项式存在，例如取一个插值多项式），都有 $\operatorname{ad}_{f(X_s)} = g(\operatorname{ad}_{X_s})$. 于是只须证明 $\operatorname{ad}_{X_s}(\mathfrak{g}) \subset \mathfrak{g}$. 而 $\operatorname{ad}_{X_s}$ 是 $\operatorname{ad}_X$ 的 Jordan 分解中的半单部分（见下面的注），故它是 $\operatorname{ad}_X$ 的多项式；既然 $\operatorname{ad}_X(\mathfrak{g}) \subset \mathfrak{g}$，证明完毕.
\end{proof}
\begin{remark*}
我们用到了如下事实：若 $X = X_s + X_n$ 是 $K$-向量空间 $V$ 上一个 $K$-线性自同态的 Jordan 分解，则 $\operatorname{ad}_X = \operatorname{ad}_{X_s} + \operatorname{ad}_{X_n}$ 是 $\mathfrak{gl}(V)$ 的自同态 $\operatorname{ad}_X$ 的 Jordan 分解. 事实上，前面已经指出 $\operatorname{ad}_{X_n}$ 幂零（Engel 定理下方的习题），而上面的证明中刚解释了 $\operatorname{ad}_{X_s}$ 可对角化. 于是只须注意 $\operatorname{ad}_{X_s}$ 与 $\operatorname{ad}_{X_n}$ 交换.
\end{remark*}
\begin{corollary*}[Cartan 判据]
设 $\mathfrak{g}$ 是一个有限维 $K$-Lie 代数. 那么 $\mathfrak{g}$ 可解当且仅当 $B_{\mathrm{ad}}(\mathfrak{g}, [\mathfrak{g}, \mathfrak{g}]) = 0$.
\end{corollary*}
\begin{proof}
$\Rightarrow$. 若 $\mathfrak{g}$ 可解，则由\newterm{ssl-lie-theorem}{Lie 定理}{Lie's theorem}，存在 $\mathfrak{g}$ 中的一个旗 $\mathcal{V} = (\mathfrak{h}_1 \subset \cdots \subset \mathfrak{h}_n)$，使 $\operatorname{ad}(\mathfrak{g}) \subset \mathfrak{b}(\mathcal{V})$. 于是 $\operatorname{ad}([\mathfrak{g}, \mathfrak{g}]) \subset [\operatorname{ad}(\mathfrak{g}), \operatorname{ad}(\mathfrak{g})] \subset \mathfrak{b}^{\mathrm{nilp}}(\mathcal{V})$，故对一切 $X \in \mathfrak{g}$ 与 $Y \in [\mathfrak{g}, \mathfrak{g}]$，有 $\operatorname{ad}_X \operatorname{ad}_Y \in \mathfrak{b}^{\mathrm{nilp}}(\mathcal{V})$，从而 $\operatorname{tr}(\operatorname{ad}_X \operatorname{ad}_Y) = 0$.
$\Leftarrow$. 设 $\mathfrak{h} := \operatorname{ad}([\mathfrak{g}, \mathfrak{g}])$，其中 $\operatorname{ad} : \mathfrak{g} \longrightarrow \mathfrak{gl}(\mathfrak{g})$ 是 $\mathfrak{g}$ 的伴随表示. 由于 $\mathfrak{h} \simeq [\mathfrak{g}, \mathfrak{g}]/(\mathfrak{z}(\mathfrak{g}) \cap [\mathfrak{g}, \mathfrak{g}])$，且 $\mathfrak{g}/[\mathfrak{g}, \mathfrak{g}]$ 是 Abel 的，可见 $\mathfrak{g}$ 可解当且仅当 $\mathfrak{h}$ 可解（见定义 2.3.1 之下习题的 ii)(b)\footnote{译注：同前注，``2.3.1''应为定义~\ref{def:ssl-2-4-1} 之误.}）. 而 $\mathfrak{h}$ 的可解性由前一命题得出.
\end{proof}
\end{numpara}
